% Pursuing Stacks (cotes 134-2 to 134-8): archive cote 134-4, site page 53 (PDF page 54), typescript page 141.
% typescript. Draft compared: ps3.tex:2689-2746.
% AI-transcribed by Claude Opus 5.5 (claude-opus-5-5, adaptive thinking, effort high), Message Batches API, batch ps2 (msgbatch_01NVwr1Yuo8qJf5rykeGBXPK), 2026-11-12; not yet checked by eye.
\scanpage{134-4/53}
\mspage{141}

totally O-connected,  we still get a description of $\hand{\Sigma}_h$ ,  as those inter-
vals such that for any \inserted[typed]{O-connected} $X$, the corresponding sections of $X\times I$ over $X$ corres-
pond to the same connected component of \xout{4x} $X\times I$.  It is enough for this
that the common  connected component $I_o$ of $I$ for $\hand{\delta}_o,\hand{\delta}_1$  should be \xout{\unsure{universal}}
O-connected over e  -  it is  not clear to me whether this condition is
equally necessary.  Anyhow,  presumably the case M totally O-connected will
be enough for all we'll have to do.

In case M is totally O-connected,  the homotopy notions in M for the
canonical homotopy structure are just those which can be described in
terms of ``homotopies'' using connected intervals  -  which is intuitively
the very first thing that comes to mind indeed,  when trying to \xout{express} \inserted[typed]{mimic}
most naïvely,in an abstract categorical context,the familiar homotopy
notions for topological spaces.

Now let
\[ W\subset Fl(M) \]
be any saturated set of arrows in M  (viewed as a notion of ``weak equiva-%
lence'' in M).  Consider the corresponding homotopy structure $h_W$ on M,
defined in terms of W-aspheric \xout{weak} intervals as a generating family of
weak homotopy intervals.  Let $h_M$\xout{\illegible{1}}  be the canonical homotopy structure
on M.  Thus the condition
\[ h_W \hand{\leq} h_M \]
just means that $h_W$ is O-admissible,  or equivalently,  that W-aspheric
objects over e which have a section,  are \xout{xxxx} O-connected\xout{\illegible{1}}  (for simpli-%
city,  I'll assume from now on M totally O-connected).  This looks like
a very reasonable condition indeed,  if W should correspond at all to the
intuitions associated to the notion of ``weak equivalence'' !  As a matter
of fact,  this condition is clearly implied by the condition that W itself
should be \xout{xxxxx}``\uline{O-admissible}'' ,  by which we mean that the functor $\hand{\pi}_o$ is
compatible with W,  i.e. transforms weak equivalences into bijections,  or
equivalently,  factors through $M\hand{\rightarrow}W^{-1}M$.

What we are looking for however is conditions on W for the \uline{equality}
\[ h_W = h_M \]
to hold.  When the previous condition \xout{is} (expressing $h_W\hand{\leq}h_M$ )  is satis-%
fied,  all that remains is to express the opposite inequality,  which is
done in the standart way.  We thus get:

\margin{53}
